% ============================================================
%  Section 3 – Method
%  Section 4 – Experimental Apparatus
%  Draft — Think Twice: Does Repeating Reasoning Summaries
%           Increase Chain-of-Thought Efficiency?
% ============================================================

% --------------------------------------------------
\section{Method}
\label{sec:method}

\subsection{Task Formulation}
\label{subsec:task}

Let $\mathcal{D} = \{(x_i, y_i^*)\}_{i=1}^{N}$ be a dataset of problems $x_i$ paired with
ground-truth answers $y_i^*$, and let $s_i$ denote a pre-computed, abstract reasoning skill
associated with $x_i$.
We define a \emph{prompt constructor} $\phi : (\mathcal{S}^k, \mathcal{X}) \to \mathcal{P}$
that assembles an input sequence from $k \geq 0$ copies of a context string and the problem:
\begin{equation}
  \phi(c, k, x) = \underbrace{c \| c \| \cdots \| c}_{k\text{ copies}} \| x,
  \label{eq:phi}
\end{equation}
where $\|$ denotes string concatenation and $c$ is one of three possible context objects (see
below).  For $k = 0$ we recover the direct prompting baseline $\phi(\cdot, 0, x) = x$.

\paragraph{Experimental conditions.}
We define four prompt conditions that are matched on copy count $k$:

\begin{itemize}
  \item \textbf{Direct (D)}: Input $= x$.  No injected context.  This is the zero-shot baseline.
  \item \textbf{Summary Repetition (\(\mathrm{SR}_k\))}: Input $= \phi(s, k, x)$.
        The pre-computed skill text $s$ is prepended $k$ times before the problem.
        At $k{=}1$ this reduces to the single-insight injection of I2S
        \cite{wang2025harm}, establishing a clean entry point.
  \item \textbf{Prompt Repetition (\(\mathrm{PR}_k\))}: Input $= \phi(x, k, x)$, i.e.\ the
        problem statement itself repeated $k{+}1$ times.
        This replicates the best-performing setting of \citet{leviathan2025prompt}
        at $k{=}1$ and extends it to higher counts.
  \item \textbf{Raw Reasoning Transfer (\(\mathrm{RT}_k\))}: Input $= \phi(\tau, k, x)$,
        where $\tau$ is a full model-generated reasoning trace for a closely related problem,
        frozen before the sweep begins.
        This instantiates the harmful few-shot baseline identified in \citet{wang2025harm}.
\end{itemize}

\paragraph{Repetition count sweep.}
We vary $k \in \{1, 2, 3, 4, 5\}$ for all injected conditions
(SR, PR, RT), yielding a dense curve of 5 points per condition.
The range is motivated as follows.
$k{=}1$ corresponds to the I2S single-copy baseline \cite{wang2025harm} and to the
standard prompt repetition of \citet{leviathan2025prompt}.
$k{=}2$ matches the primary condition reported by \citeauthor{leviathan2025prompt},
who found it optimal for non-reasoning models on most benchmarks.
$k{=}3$ was also evaluated by \citeauthor{leviathan2025prompt} and shown to
occasionally outperform $k{=}2$, particularly on tasks with high positional sensitivity.
$k \in \{4, 5\}$ extends beyond the previously studied range to probe
the \emph{cognitive overload} regime, where extraneous context is expected
to exceed the model's effective working memory and degrade performance,
consistent with Cognitive Load Theory as applied to transformer positional encodings
in recent architectural analyses.
Together, the five points map the full accuracy-versus-repetition curve and allow us to
identify whether a performance peak exists before degradation sets in.

\paragraph{Abstraction levels of the skill text.}
The content of the injected skill $s$ is not uniform: following the taxonomy of
\citet{didolkar2025metacognitive}, who distinguish behaviors by their level of
procedural specificity, we test three abstraction levels drawn from the skill texts
available in our dataset:

\begin{enumerate}
  \item \textbf{Strategy Steps (SR-Steps)}: A numbered sequence of procedural steps
        specific to the problem type (e.g.\ ``Step 1: identify the variable.
        Step 2: isolate it algebraically\ldots'').
        This is the richest, most concrete level.
  \item \textbf{Meta-Cognitive Hint (SR-Hint)}: A single-sentence strategic pointer
        without procedural detail (e.g.\ ``Use proportional reasoning to relate the
        given quantities''). This matches the concise \emph{behavior name + instruction}
        format of \citet{didolkar2025metacognitive}.
  \item \textbf{Structure Template (SR-Template)}: An abstract schema of the solution
        form with placeholder variables (e.g.\ ``If $A$ relates to $B$ by ratio $r$,
        then $C = r \cdot B$'').
        This is the most abstract level, testing whether structural pattern transfer
        alone is sufficient.
\end{enumerate}

Comparing these levels allows us to separate \emph{how much semantic content} the
repeated object must carry to be effective, independently of the repetition count $k$.

\subsection{Architecture Design \& Context Positioning}
\label{subsec:architecture}

\paragraph{Prompt injection.}
All injected context is inserted as part of the user turn in the model's chat template,
before the instruction to answer the question.
No system-prompt manipulation is performed; the model's default system prompt is retained
across all conditions to avoid confounding positional effects with system-prompt effects.

\paragraph{Context position variants.}
The position of the repeated context within the prompt is itself an experimental variable.
We test three placements:

\begin{enumerate}
  \item \textbf{Pre-problem} (primary): $[\mathrm{context}^k]\ [x]$ — the skill text or
        repeated query precedes the problem statement.
        Based on the finding that causal attention in autoregressive models is
        asymmetric — later tokens attend to earlier ones but not vice versa — placing
        context early ensures maximum downstream attention coverage.
  \item \textbf{Post-problem}: $[x]\ [\mathrm{context}^k]$ — the context follows the
        problem.  This tests whether summary information placed after the problem
        still influences generation, and is used as a control for the \emph{Lost in
        the Middle} effect \cite{liu2023lost}, which predicts a performance dip for
        information placed in the center of a long sequence.
  \item \textbf{Within thinking tokens} (reasoning models only): the skill text is
        injected as a preamble inside the explicit \texttt{<think>} block provided
        to the model, simulating a \emph{warm-started} reasoning trace.
        This placement is only applicable in the thinking-enabled mode of Qwen3-4B
        and tests whether procedural context inserted directly into the reasoning
        scratchpad reduces subsequent deliberation tokens more effectively than
        prompt-level injection.
\end{enumerate}

The pre-problem placement is used as the default for all main results; the post-problem
and within-thinking variants are evaluated in the ablation study
(Section~\ref{subsec:ablation}).

\subsection{Training \& Evaluation Objective}
\label{subsec:objective}

No model parameters are updated at any point.
All interventions are purely prompt-level.
We define a simple evaluation objective over a held-out subset of $\mathcal{D}$:
\begin{equation}
  \mathcal{L} = \mathbb{E}_{x \sim \mathcal{D}} \left[ \ell\!\left(f_\theta(x), y^*\right) \right],
  \label{eq:objective}
\end{equation}
where $f_\theta$ is the frozen language model, $x$ is the fully assembled prompt for a
given condition and copy count, and $\ell$ is the task-specific loss (exact-match accuracy
for math problems; see Section~\ref{subsec:metrics} for details).
Token efficiency is measured as a secondary objective: for each inference call we record
the number of tokens generated inside the \texttt{<think>} block and the total output
length, providing a cost-accuracy trade-off curve as a function of $k$.

% --------------------------------------------------
\section{Experimental Apparatus}
\label{sec:apparatus}

\subsection{Datasets}
\label{subsec:datasets}

\paragraph{Primary dataset — TRS Reasoning-Skill.}
Our primary dataset is the \texttt{stallone0000/Reasoning-Skill} collection on HuggingFace,
which provides (problem, answer, skill\_text) triples for mathematical reasoning tasks.
This dataset was specifically chosen to \emph{isolate the repetition effect from the
insight-acquisition effect}: because skill texts are pre-computed and bundled with each
problem, no auxiliary retriever or skill-extraction model is introduced.
Any such auxiliary component would act as a confound, entangling the quality of the
extracted insight with the effect of its repetition.
The math domain also provides objective, binary evaluation (exact-match correctness)
and aligns with the benchmarks used in related work
\cite{leviathan2025prompt, wang2025harm, didolkar2025metacognitive}.

\paragraph{Secondary datasets — generalization benchmarks.}
To assess whether findings generalise beyond mathematical reasoning, we additionally
evaluate on three complementary benchmarks:

\begin{itemize}
  \item \textbf{GSM8K} \cite{cobbe2021gsm8k} / \textbf{MATH} \cite{hendrycks2021math}:
        Standard mathematical reasoning benchmarks used in both
        \citet{leviathan2025prompt} and \citet{wang2025harm}, enabling direct
        comparison with reported results.
  \item \textbf{BBH — Penguins in a Table} and \textbf{BBH — Disambiguation QA}
        \cite{suzgun2022bbh}: Two BIG-Bench Hard subsets that probe tabular
        reasoning and reference disambiguation respectively.
        These tasks are chosen because they require distinct reasoning strategies,
        testing whether the skill-repetition effect transfers across problem types.
  \item \textbf{ProofWriter} \cite{tafjord2021proofwriter}: A logical deduction
        benchmark in which skill texts (inference rules) are naturally structured
        and available, making it well-suited to the SR-Template abstraction level.
\end{itemize}

\subsection{Preprocessing}
\label{subsec:preprocessing}

\paragraph{Skill text extraction.}
For each problem in the evaluation split we retrieve the corresponding
\texttt{skill\_text} field from the TRS dataset and use it directly as $s$ in
the SR conditions.
No further processing is applied to the TRS skill texts, preserving the three
abstraction levels described in Section~\ref{subsec:task} as they appear in
the dataset.

\paragraph{Raw reasoning trace generation.}
For the RT conditions, we generate one reasoning trace per problem by querying
Qwen3-4B with thinking enabled and no injected context.
The full content of the \texttt{<think>}\ldots\texttt{</think>} block is extracted
programmatically and stored as $\tau$.
These traces are generated once and frozen before any experimental sweep begins,
ensuring that RT traces are not confounded by the model's own adaptation over the course
of the experiment.

\paragraph{Tokenisation and length accounting.}
All token counts are obtained from the model's own generation metadata
(thinking tokens and answer tokens are reported separately by the Qwen3-4B API).
We report \emph{total generated tokens} (thinking + answer) as the primary efficiency
metric, since this reflects the actual inference cost, not only the visible output
\cite{leviathan2025prompt}.

\subsection{Procedure \& Baselines}
\label{subsec:procedure}

\paragraph{Model variants.}
Our primary controlled comparison uses a single model in two operating modes:
\begin{itemize}
  \item \textbf{Qwen3-4B / thinking on} (\textsc{Qwen-Think}):
        Reasoning model; generates an explicit chain of thought inside
        \texttt{<think>} tags before the final answer.
  \item \textbf{Qwen3-4B / thinking off} (\textsc{Qwen-NoThink}):
        Non-reasoning model; produces a direct answer without a scratchpad.
        This mirrors the non-reasoning setting in \citet{leviathan2025prompt},
        where prompt repetition was found most beneficial.
\end{itemize}
Using the same weights in both modes eliminates model capacity as a confound between
reasoning and non-reasoning comparisons.
A \textbf{second model family} (TBD; candidate: Gemma~3 or Llama~3) is included
to distinguish Qwen-specific behavior from more general effects, as called for in the
contributions of the paper.

\paragraph{Experimental conditions matrix.}
For each dataset, each model variant, and each condition
$c \in \{\mathrm{D}, \mathrm{SR}_k, \mathrm{PR}_k, \mathrm{RT}_k\}$
with $k \in \{1, 2, 3, 4, 5\}$, we run the model on all problems in the evaluation split.
The full matrix is:
$2\ \text{model modes} \times 5\ \text{conditions} \times 5\ \text{copy counts}
= 50\ \text{configurations}$,
plus the single Direct baseline per model mode.

\paragraph{Noise ablation baselines.}
To isolate whether performance changes stem from the \emph{content} of the injected
text rather than its mere token-length increase, we include two length-matched noise
conditions at each $k$:

\begin{itemize}
  \item \textbf{Filler ($\mathrm{Fill}_k$)}: The context is replaced by a sequence
        of period characters padded to the same token count as $\mathrm{SR}_k$.
        This replicates the \emph{Padding} control of \citet{leviathan2025prompt},
        who confirmed that padding does not improve performance, and is consistent
        with the mechanistic prediction of \citet{shi2025meaningless} that even
        meaningless tokens can induce activation redistribution.
        Any improvement of SR over Fill is therefore attributable to semantic content.
  \item \textbf{Irrelevant Reasoning ($\mathrm{IRel}_k$)}: The context is a
        reasoning trace drawn from a problem in a \emph{different domain}
        (e.g.\ a logical-deduction trace injected before a math problem),
        token-count matched to $\mathrm{RT}_k$.
        This tests whether domain-matched content is necessary, or whether any
        plausible-looking reasoning text provides the same benefit.
\end{itemize}

\paragraph{Statistical validity.}
All conditions are run with 3 independent random seeds (controlling prompt-order
shuffling and any stochastic decoding). We report mean accuracy and mean token counts
together with standard deviations. Statistical significance between conditions at
matched $k$ is assessed with a paired McNemar test \cite{mcnemar1947} at
$p < 0.1$, following \citet{leviathan2025prompt}.

\subsection{Metrics}
\label{subsec:metrics}

\begin{itemize}
  \item \textbf{Accuracy}: For math tasks (TRS, GSM8K, MATH), exact-match
        correctness of the final numerical answer, extracted by stripping formatting.
        For BBH and ProofWriter, we follow the standard evaluation protocol of
        each benchmark \cite{suzgun2022bbh, tafjord2021proofwriter}.
  \item \textbf{Thinking token count} ($T_\mathrm{think}$): The number of tokens
        generated inside the \texttt{<think>} block (applicable only to
        \textsc{Qwen-Think}).
        A reduction in $T_\mathrm{think}$ signals that the injected context
        displaced some deliberation, consistent with the token-saving hypothesis.
  \item \textbf{Total generated token count} ($T_\mathrm{total}$): Thinking tokens
        plus answer tokens.  This is the primary efficiency metric because it
        captures the full inference cost, including hidden reasoning, rather than
        only the visible output \cite{leviathan2025prompt}.
  \item \textbf{Accuracy–efficiency trade-off}: Accuracy plotted against
        $T_\mathrm{total}$ across values of $k$ for each condition,
        to identify whether any configuration simultaneously improves accuracy
        and reduces generated tokens.
\end{itemize}

\subsection{Main Results}
\label{subsec:main-results}

% TODO: populate after experiments are complete.
% Key comparisons to report:
%   SR_k vs. D       → net effect of injected summary repetition
%   SR_k vs. PR_k    → semantic value of summary over mechanical repetition
%   SR_k vs. RT_k    → value of distillation (abstract vs. raw trace)
%   SR_k(k=1) vs. SR_k(k>1) → pure stacking effect within summary condition
%   Qwen-Think vs. Qwen-NoThink → reasoning vs. non-reasoning mode difference

\begin{table*}[t]
  \caption{%
    Accuracy (\%) and total generated token count ($T_\mathrm{total}$, in parentheses)
    for all conditions across repetition counts $k \in \{1,\ldots,5\}$.
    \textsc{Qwen-Think} uses thinking enabled; \textsc{Qwen-NoThink} uses thinking
    disabled.
    \textbf{Bold} marks the best accuracy per model mode.
    $\dagger$ denotes a statistically significant improvement over Direct~(D)
    at $p < 0.1$ (McNemar test).
    [-- ] indicates the condition is not applicable (e.g.\ thinking token counts for
    the non-reasoning model).
  }
  \label{tab:main-results}
  \vskip 0.15in
  \begin{center}
  \begin{small}
  \begin{tabular}{lcccccc}
    \toprule
    & \multicolumn{3}{c}{\textit{\textsc{Qwen-Think} (reasoning on)}}
    & \multicolumn{3}{c}{\textit{\textsc{Qwen-NoThink} (reasoning off)}} \\
    \cmidrule(lr){2-4} \cmidrule(lr){5-7}
    Condition ($k$)
      & Acc.\ (\%) & $T_\mathrm{think}$ & $T_\mathrm{total}$
      & Acc.\ (\%) & $T_\mathrm{think}$ & $T_\mathrm{total}$ \\
    \midrule
    \multicolumn{7}{c}{\textit{TRS Reasoning-Skill (math)}} \\
    Direct (D)            & --   & --  & --  & --   & {[--]} & --  \\
    $\mathrm{SR}_1$       & --   & --  & --  & --   & {[--]} & --  \\
    $\mathrm{SR}_2$       & --   & --  & --  & --   & {[--]} & --  \\
    $\mathrm{SR}_3$       & --   & --  & --  & --   & {[--]} & --  \\
    $\mathrm{SR}_4$       & --   & --  & --  & --   & {[--]} & --  \\
    $\mathrm{SR}_5$       & --   & --  & --  & --   & {[--]} & --  \\
    \midrule
    $\mathrm{PR}_1$       & --   & --  & --  & --   & {[--]} & --  \\
    $\mathrm{PR}_3$       & --   & --  & --  & --   & {[--]} & --  \\
    $\mathrm{PR}_5$       & --   & --  & --  & --   & {[--]} & --  \\
    \midrule
    $\mathrm{RT}_1$       & --   & --  & --  & --   & {[--]} & --  \\
    $\mathrm{RT}_3$       & --   & --  & --  & --   & {[--]} & --  \\
    $\mathrm{RT}_5$       & --   & --  & --  & --   & {[--]} & --  \\
    \midrule
    $\mathrm{Fill}_1$     & --   & --  & --  & --   & {[--]} & --  \\
    $\mathrm{Fill}_3$     & --   & --  & --  & --   & {[--]} & --  \\
    $\mathrm{Fill}_5$     & --   & --  & --  & --   & {[--]} & --  \\
    \bottomrule
  \end{tabular}
  \end{small}
  \end{center}
  \vskip -0.1in
\end{table*}

\subsection{Ablation Study}
\label{subsec:ablation}

% TODO: populate after ablation experiments are complete.
% Two ablation axes:
%   1. Context positioning (pre-problem / post-problem / within thinking)
%   2. Noise baselines (Filler / Irrelevant Reasoning) at matched token length

\begin{figure}[h]
  \centering
  \includegraphics[width=\linewidth]{example-image}
  \caption{%
    Accuracy as a function of repetition count $k$ for all four conditions
    (\textsc{D}, $\mathrm{SR}_k$, $\mathrm{PR}_k$, $\mathrm{RT}_k$) under
    \textsc{Qwen-Think} (solid) and \textsc{Qwen-NoThink} (dashed).
    The shaded band shows $\pm 1$ standard deviation across seeds.
    A peak followed by decline would support the cognitive-overload hypothesis.
  }
  \label{fig:accuracy-vs-k}
\end{figure}

\subsection{Qualitative Results}
\label{subsec:qualitative}

% TODO: include 2–3 representative examples after experiments are complete.
% Each example should show side-by-side outputs for:
%   Direct (D) | SR_1 | SR_3 | SR_5
% highlighting changes in reasoning verbosity and error patterns.
